\documentclass[10pt]{article}
\usepackage[letterpaper,margin=0.60in]{geometry}
\usepackage{amsmath,amssymb,url}
\pagestyle{empty}
\setlength{\parindent}{0pt}
\setlength{\parskip}{3pt}
\setlength{\abovedisplayskip}{4pt}
\setlength{\belowdisplayskip}{4pt}
\begin{document}
\small
\begin{center}
\textbf{The finite-circle conformal cut and its lattice domain}
\end{center}
\vspace{-5pt}
\textbf{Scope [R].}
This derivation serves the Sec.~V comparator and control D1.  A massless
sector with no centrifugal term, a reflecting centre and outer wall,
unfolds to a chiral complex fermion, \(c_{\rm CFT}=1\), on an initial
optical circle of circumference \(L\).  Assume its pure Neveu--Schwarz
vacuum, \(\psi(p+L)=-\psi(p)\), with unchanged reflection phases.
Neither regularity at the centre nor absence of a zero level alone
certifies this monodromy.  There are two sectors per chain.

For \(ds^2=-N^2dt^2+(dR-v\,dt)^2\), let \(p_o,p_i\) be outgoing
and incoming initial-vacuum labels at a fixed exterior cut \(R_c\).
Dots denote fixed-cut time derivatives;
\(\dot p_o,\dot p_i>0\), \(D=N^2-v^2>0\),
and \(0<\ell=p_i-p_o<L\).  Put
\(d_L=(L/\pi)\sin(\pi\ell/L)\), with all reflections included.
Replica uniformization gives \(h_n=(n-1/n)/24\); the circle
two-point function yields
\(S=\frac16\ln[d_L/\sqrt{\epsilon_o\epsilon_i}]+C\).
The decisive endpoint step, for lattice spacing \(a\), follows from
\((\partial_t+(v\pm N)\partial_R)p_{o,i}=0\):
\[
 \epsilon_o=\frac{a\dot p_o}{N+v},\qquad
 \epsilon_i=\frac{a\dot p_i}{N-v},\qquad
 S_{\rm conf}=\frac16\ln\frac{d_L}a
       -\frac1{12}\ln(\dot p_o\dot p_i)+\frac1{12}\ln D+C .
 \tag{1}
\]
Purity makes this the exterior entropy.  Subtract the same initial
vacuum, denoted by \(0\), with the same cutoff convention, so \(C\)
cancels.  Set \(F=\frac16\ln(d_L/a)\),
\(n_j=-\ln(\dot p_j/\dot p_{j0})\), and
\(A=[n_i+\ln(D/D_0)]/12\).  The e-folding identity and its rate
on an exponential segment \(\dot n_o=\kappa\) are
\[
 B_{\rm conf}\equiv\Delta S_{\rm conf}-\Delta F=\frac{n_o}{12}+A,
 \qquad
 \dot B_{\rm conf}=\frac{\kappa}{12}
       +\frac1{12}\frac{d}{dt}\ln\frac{D}{\dot p_i}.
 \tag{2}
\]
Thus \(\dot B_{\rm conf}=\kappa/12\) precisely when \(D/\dot p_i\)
is constant.  Raw \(S_{\rm conf}\) adds
\(\dot F=\frac{\pi}{6L}\cot(\pi\ell/L)(\dot p_i-\dot p_o)\);
its rate is \(\kappa/12\) when both extra terms sum to zero.
For a supplied map \(p(u)\), use
\(\dot p_o=p'(u_c)\dot u_c\), including this clock factor in \(n_o\).
The reference here is the initial vacuum.

Equation (1) is a conformal scaling formula.  A finite chain can agree
when endpoint separations are resolved, fields and label maps are smooth
on the lattice scale, and accumulated propagation errors are small.
For the local constant-field symbol, \(k\) is displacement from a massless
valley and \(s=\pm1\) includes its chirality:
\[
 \omega_s=(v+sN)k-\frac{a^2}{6}(4v+sN)k^3+O(a^4k^5),\qquad
 a|k|\ll1,\quad
 \frac{a^2|4v+sN|k^2}{6|v+sN|}\ll1 .
 \tag{3}
\]
These conditions must hold along the relevant rays.  Precursors
blueshift as \(e^{n_o}\); cancellation in \(v+N\) tightens (3).
The finite-spacing parting step is replacing the finite-band propagator
by conformal transport outside this domain.

The source-reported checks, in nats per sector, remain numerical targets.
The 1200-site prescribed-flow harness gives \(7.1\times10^{-5}\)
through two e-foldings and \(2.1\times10^{-4}\) through four,
with its first \(10^{-3}\) departure at 4.07.
For B2 use the fixed unit \(r_0=64a\): at \(R_c=2r_0\), the reported
error is \(4\times10^{-4}\) through about two e-foldings and
\(6.3\times10^{-3}\) (A2), \(6.4\times10^{-3}\) (D1), at 2.90.
Departures above 0.01 begin at 2.85--3.47 across five cuts.
D1's separate \(7\times10^{-5}\) is tracer step refinement.
There is no universal three-e-folding boundary; the incoming-label
correction for \(p_i>0\), quoted as at most \(2\times10^{-3}\), is estimated.

Define the lattice observable \(B_{\rm lat}=\Delta S_{\rm lat}-\Delta F\);
the Jacobian equality (2) is restricted to the conformal description.
At the fixed \(2r_0\) cut, the supplied \(0.92\pm0.04\) target on
\(60\le t/r_0\le100\) uses A2's growth-minus-circle convention
(reported mean 0.925).  D1's header additionally subtracts \(A\);
its similarly rounded number describes a different observable.
For \(B_{\rm lat}\), compare the supplied cutoff ansatz
\(B_{\rm sat}(c)=\ln[c/(aT_H)]/12+\delta\), identified as withdrawn
in A2, with \(c\) distinct from \(c_{\rm CFT}\):
\[
 r_E(t)=2GE(t),\qquad T_H(t)=\frac1{4\pi r_E(t)},\qquad
 B_\pi(t)=\frac1{12}\ln\!\left[4\pi^2\frac{r_E(t)}a\right]+\delta .
 \tag{4}
\]
The charge radius \(r_E\) is distinct from the applied trapping root.
At \(t=100r_0\), supplied entries give \(r_E/a=53.03\) and
\(B_{\rm lat}=1.3872-0.4453=0.9419\), against \(B_\pi\simeq0.637\)
for \(\delta=0\), or 0.633 for the separate \(\delta=-0.004\) convention.
The value 0.653 uses the nominal \(r_E=r_0\).
These are same-clock arithmetic comparisons of source entries;
a window comparison needs the same-clock, same-weight average of (4),
which is not recomputed here.  An outgoing-only interpretation of the
ansatz needs \(A\) added for comparison with \(B_{\rm lat}\).
No numerical result is independently cleared and no saturation constant
is derived.

\vspace{2pt}
{\footnotesize\textbf{SOURCES.} The independent draft used isolated
formula/map statements and supplied definitions; full A2 and D1 text
was read for the review corrections.\\
\path{<internal>/REVIEWA_write_A2_v1.1.md}
\\
\path{<internal>/PAGE_VERIFY_D1_conformal_control_v1.1.md}
\\
\path{<internal>/redteam_CONFORMAL_cut_formula_second_route_v1.md}
\\
\path{<internal>/requirements_new.txt}.}
\end{document}
